\section{The signed-service seam: a full physical-time
proof}

\subsection{1. Scope and result}

Write \(D=B+C\) for the mixed complex and

\[
 W=B-C.
\tag{1.1}
\]

This note treats a pure-\(C\) linkage \(L_C\) with support

\[
 \{C,2C\}
 \quad\hbox{or}\quad
 \{0,C,2C\},
\tag{1.2}
\]

and a disjoint mixed linkage \(L_M\) whose support is one of

\[
 \{0,A,D\},\qquad
 \{0,2A,D\},\qquad
 \{A,2A,D\},\qquad
 \{0,A,2A,D\}.
\tag{1.3}
\]

Every labelled reaction graph is arbitrary subject to strong
connectivity, and all rate constants are positive. Parallel channels are
aggregated only to simplify notation. Stochastic propensities are
falling-factorial mass-action propensities. The \(A\leftrightarrow B\)
images are obtained by relabeling.

The first support in (1.2) includes the exceptional architecture

\[
 \{C,2C\}\ \&\ \{0,A,2A,B+C\},
\tag{1.4}
\]

and the second includes

\[
 \{0,C,2C\}\ \&\ \{A,2A,B+C\}.
\tag{1.5}
\]

The conclusion is classwise.

\begin{quote}
\textbf{Signed-service seam theorem.} Every closed irreducible
population class of every network in (1.2)--(1.3) is positive recurrent.
The process is nonexplosive. No recurrence assertion for the raw
embedded jump chain is used.
\end{quote}

The proof has two parts. A proper quadratic potential gives ordinary
physical-time generator drift whenever \(A\) or \(C\) is large, or
whenever \(BC\) is large. The only remaining tube has bounded \(A\),
\(C=0\), and \(B\to\infty\). On that tube a fast-phase return cycle has
a \(B\)-decreasing service probability of order \(1/B\), an
adverse-service probability of order \(1/B^2\), and uniformly bounded
physical duration. The quadratic potential converts this into a uniform
negative return-cycle drift.

\subsection{\texorpdfstring{2. Why \(V+\gamma(B-C)\) is not the
repair}{2. Why V+\textbackslash gamma(B-C) is not the repair}}

The linear signed correction is proper when added to a factorial
potential, but it is too small at the service bottleneck.

Consider the strongly connected mixed cycle

\[
 0\longrightarrow D\longrightarrow2A\longrightarrow A\longrightarrow0
\tag{2.1}
\]

together with \(C\rightleftarrows2C\). Start at

\[
 x_N=(0,N,0).
\tag{2.2}
\]

The first mixed reaction \(0\to D\) produces \((0,N+1,1)\). At that
target, the probability that \(C\to2C\) fires before the \(D\)-source
reaction is

\[
 \frac{\alpha}{d(N+1)+O(1)}=\Theta(N^{-1}),
\tag{2.3}
\]

where \(\alpha\) is the \(C\to2C\) rate constant and \(d\) is the
aggregate rate constant out of \(D\). On failure, the mixed cycle
returns to the bounded phase in \(O(1)\) mean time. On success, two
rapid \(D\)-source clearings, followed by return of the \(A\)-phase,
take \(x_N\) to \(x_{N-1}\) with probability \(1-O(N^{-1})\).

Thus a bounded-mean regeneration cycle has service probability
\(\Theta(N^{-1})\). For

\[
 F(x)=\sum_i\log(x_i!)+\gamma(B-C),
\tag{2.4}
\]

the successful return reward is only

\[
 F(x_{N-1})-F(x_N)=-\log N-\gamma.
\tag{2.5}
\]

Consequently the negative expected reward per regeneration cycle is only
\(O((\log N)/N)\), while the mean physical duration stays bounded away
from zero. No fixed \(\eta>0\) can give

\[
 \mathbb E[\Delta F+\eta\tau]\le-\varepsilon
\tag{2.6}
\]

through these cycles. Pointwise drift fails even more plainly: at
\(x_N\), the reaction \(0\to D\) has factorial reward \(\log(N+1)\), and
the signed term is unchanged.

This does not disprove recurrence. It shows that the proposed
factorial-plus-linear signed repair has the wrong growth. A potential
\(f(B)\) needs

\[
 f(N)-f(N-1)=\Omega(N)
\tag{2.7}
\]

to convert a service probability of order \(1/N\) into constant
return-cycle drift. Quadratic growth is the minimal natural choice.

\subsection{3. A common product-form shortcut for three mixed
complexes}

When \(L_C=\{C,2C\}\) and \(L_M\) is one of the first three supports in
(1.3), both linkages have deficiency zero. This already closes those
three cases.

Indeed, \(C\rightleftarrows2C\) has complex-balanced value

\[
 \theta_C=\alpha/\beta.
\tag{3.1}
\]

The mixed linkage has invariant \(B-C\). Its complex-balance equations
determine \(\theta_A\) and the product \(\theta_B\theta_C\), but the
invariant scaling

\[
 (\theta_B,\theta_C)\mapsto(e^s\theta_B,e^{-s}\theta_C)
\tag{3.2}
\]

leaves every mixed monomial unchanged. Hence its complex-balanced vector
can always be chosen with the value (3.1). The two linkage balances then
hold at one common vector. The product-Poisson measure at that vector,
restricted to any irreducible population class, is a summable invariant
measure. This proves positive recurrence on every closed class.

The four-complex support in (1.3) has deficiency one, so this shortcut
is not available for arbitrary rates. The physical-time proof below
treats it and also gives a uniform proof for all supports in (1.3),
including the second pure-\(C\) support in (1.2).

\subsection{4. A rate-dependent positive phase
workload}

Let

\[
 u(0)=0,\qquad u(A)=p_A,\qquad u(2A)=2p_A,
 \qquad u(D)=p_B+p_C,
\tag{4.1}
\]

where \(p_A,p_B,p_C>0\). Put

\[
 U(a,b,c)=p_Aa+p_Bb+p_Cc.
\tag{4.2}
\]

The weights can be selected so that every highest-degree mixed source
has strictly negative \(U\)-drift.

\subsubsection{Lemma 4.1 (weight
interval)}

For every strongly connected orientation and every positive rate vector
on one of the supports (1.3), positive weights in (4.2) can be chosen
with the following properties.

\begin{itemize}
\tightlist
\item
  If \(2A\) is present, the aggregate \(U\)-increments out of both
  \(2A\) and \(D\) are strictly negative.
\item
  If \(2A\) is absent, the aggregate \(U\)-increments out of both \(A\)
  and \(D\) are strictly negative.
\end{itemize}

To see this, first suppose \(2A\) is present. Normalize

\[
 u(2A)=1,\qquad u(A)=\tfrac12,\qquad u(D)=r.
\tag{4.3}
\]

Write \(\kappa_{ij}=0\) for an absent edge. The two required
inequalities are

\[
 r<L_2^{\rm up}:=
 1+\frac{\kappa_{2,0}+\kappa_{2,A}/2}{\kappa_{2,D}}
\tag{4.4}
\]

with \(L_2^{\rm up}=+\infty\) when \(\kappa_{2,D}=0\), and

\[
 r>L_D^{\rm low}:=
 \frac{\kappa_{D,A}/2+\kappa_{D,2}}
 {\kappa_{D,0}+\kappa_{D,A}+\kappa_{D,2}}.
\tag{4.5}
\]

Terms for vertices absent from the support are omitted. One has
\(L_D^{\rm low}\le1\le L_2^{\rm up}\). Equality on both sides would say
that \(2A\) and \(D\) only point to each other, contradicting strong
connectivity with the remaining vertex or vertices. Thus the open
interval is nonempty.

For \(\{0,A,D\}\), normalize \(u(A)=1,u(D)=r\). The interval is

\[
 \frac{\kappa_{D,A}}{\kappa_{D,0}+\kappa_{D,A}}
 <r<
 1+\frac{\kappa_{A,0}}{\kappa_{A,D}},
\tag{4.6}
\]

with the same convention for a missing denominator. Strong connectivity
again makes the interval nonempty. Finally split the selected value
\(u(D)\) as \(p_B+p_C\) with both summands positive.

\subsubsection{Lemma 4.2 (phase generator
drift)}

For the weights from Lemma 4.1 there are constants \(K,k>0\) such that

\[
 \mathcal LU\le K-k\bigl(r_A(A)+BC+C^2\bigr),
\tag{4.7}
\]

where \(r_A(A)=A\) for the support \(\{0,A,D\}\), and \(r_A(A)=A^2\) for
the other three supports.

For the mixed linkage this is immediate after collecting propensities by
source complex. The \(0\)-source contribution is constant, the
\(A\)-source contribution is at most linear, and Lemma 4.1 makes the
\(2A\)- and \(D\)-source coefficients strictly negative. If \(2A\) is
absent, the \(A\)-source coefficient itself is strictly negative.

For \(L_C=\{C,2C\}\),

\[
 \mathcal L_CU=p_C[\alpha C-\beta C(C-1)].
\tag{4.8}
\]

For \(L_C=\{0,C,2C\}\), every edge out of \(2C\) lowers \(C\), so the
quadratic coefficient is again strictly negative; the \(0\)- and
\(C\)-source terms are at most constant and linear. Absorbing those
lower orders proves (4.7).

All jumps of \(U\) are bounded. The total reaction rate is at most a
constant times

\[
 1+r_A(A)+BC+C^2.
\tag{4.9}
\]

Consequently

\[
 \mathcal L(U^2)
 =2U\mathcal LU+
   \sum_r\lambda_r(x)(\Delta_rU)^2.
\tag{4.10}
\]

Equations (4.7)--(4.10) show that there are \(R<\infty\) and
\(\varepsilon>0\) such that

\[
 \mathcal L(U^2)\le-1-\varepsilon
\tag{4.11}
\]

outside a finite set and the tube

\[
 \mathcal T_R=\{(a,b,0):a\le R\}.
\tag{4.12}
\]

Indeed, bounded \(r_A(A)+BC+C^2\) forces \(A,C\) to be bounded; if
\(C\ge1\), it also forces \(B\) to be bounded. Thus (4.12) is the only
infinite region not handled by pointwise generator drift. Notice that

\[
 \Phi(x):=1+U(x)^2
\tag{4.13}
\]

is proper.

\subsection{5. The fast trace on a large positive signed
shell}

This section supplies the uniform statement that was absent from the old
finite-shell service argument.

Fix \(W=N>0\). Until a pure-\(C\) reaction occurs,

\[
 B=N+C.
\tag{5.1}
\]

Let \(d_M>0\) be the aggregate rate constant out of \(D\). At \(C\ge1\),
the aggregate mixed \(D\)-source rate is exactly

\[
 d_M(N+C)C.
\tag{5.2}
\]

Delete \(D\) from the mixed complex graph. An edge \(i\to D\), followed
by the first \(D\to j\), induces the effective rate

\[
 \bar\kappa_{ij}\supset
 \kappa_{iD}\frac{\kappa_{Dj}}{d_M}.
\tag{5.3}
\]

Direct lower-complex edges retain their rates. The resulting graph on

\[
 \{0,A\},\quad\{0,2A\},\quad\{A,2A\},
 \quad\hbox{or}\quad\{0,A,2A\}
\tag{5.4}
\]

is strongly connected: replace every occurrence of the single deleted
vertex in a simple directed path by (5.3).

The population chain of this effective one-species network is positive
recurrent on each class. In the four cases it is respectively an
immigration-death chain, a parity-preserving
pair-immigration/quadratic-death chain, a positive-population
linear-birth/quadratic-death chain, or an open one-species binary chain.
A power of \(1+A\) has negative generator drift of one degree higher
outside a finite set. In particular, return times to the minimal atom

\[
 a_*=0,\quad a_*=0\hbox{ or }1,\quad a_*=1,\quad a_*=0
\tag{5.5}
\]

for the four cases in (5.4) have finite moments. The choice \(0\) or
\(1\) in the second case is the conserved \(A\)-parity.

\subsubsection{Lemma 5.1 (uniform trace and occupation
estimates)}

Start at

\[
 e_N=(a_*,N,0)
\tag{5.6}
\]

and let \(\sigma_N\) be the next return to the same minimal
\(A,C\)-phase, after at least one reaction. First suppress pure-\(C\)
reactions. Then there are constants \(K<\infty\), \(j>0\), independent
of \(N\), such that

\[
 \sup_{N\ge N_0}\mathbb E_{e_N}\sigma_N^2\le K,
\tag{5.7}
\]

the expected number \(J_N\) of entrances into \(D\) during the cycle
obeys

\[
 j\le\mathbb EJ_N\le K,
\tag{5.8}
\]

and

\[
 \mathbb E\int_0^{\sigma_N}C_s\,ds=\Theta(N^{-1}),
 \qquad
 \mathbb E\int_0^{\sigma_N}(C_s)_2\,ds=O(N^{-2}).
\tag{5.9}
\]

The same estimates hold from any \(A\) in a fixed finite set before the
first hit of \(a_*\).

Here is a proof, including the tail point needed for uniformity. On a
fixed \(W=N\) shell, subtract the constant \(p_BN\) from \(U\) and put

\[
 R_N=p_AA+(p_B+p_C)C.
\tag{5.10}
\]

The mixed generator calculation from Lemma 4.2 gives

\[
 \mathcal L_MR_N
 \le K-k\{r_A(A)+NC+C^2\}.
\tag{5.11}
\]

Here is the promised uniform high-power estimate. Put

\[
 Q_N(A,C)=r_A(A)+NC+C^2.
\]

Every jump of \(R_N\) is bounded, and the total mixed rate is at most
\(K(1+Q_N)\). Taylor's formula for an integer \(m\ge4\), followed by
(5.11), therefore gives, outside a phase set independent of \(N\),

\[
\begin{split}
 \mathcal L_M(1+R_N)^m
 &\le m(1+R_N)^{m-1}\mathcal L_MR_N\\
 &\qquad+K_m(1+Q_N)(1+R_N)^{m-2}\\
 &\le K_m-c_m Q_N(1+R_N)^{m-1}\\
 &\le K_m-c'_m(1+R_N)^m .
\end{split}
\tag{5.11a}
\]

The last inequality uses \(1+Q_N\ge c(1+R_N)\), uniformly for
\(N\ge N_0\). Stopped Dynkin estimates, first for (5.11a) and then for
its lower powers, give uniform moments of \(A,C\) and of the time to a
fixed bounded phase set. The same stopped estimates, applied with
successively higher \(m\), control the integrated lower-source
intensities and their required moments.

Inside the bounded phase set, (5.2) makes an induced \(D\)-excursion
last \(O(N^{-1})\), and its exit target has the state-independent
probabilities \(\kappa_{Dj}/d_M\). Each edge of a fixed effective path
is therefore realized, in a uniformly bounded time, with probability
bounded away from zero. The strong effective graph (5.3) gives a uniform
minorization to the atom (5.5). Geometric repetition proves (5.7) and
the upper bound in (5.8). A fixed feasible reaction sequence from the
atom that enters \(D\) has probability bounded away from zero before
return, proving the lower bound.

Now apply the stopped estimates over this atom cycle. Because its two
endpoints have the same value of \(R_N\), (5.7), (5.11a), and the
corresponding lower-power inequalities give

\[
 \mathbb E_{e_N}\int_0^{\sigma_N}
 Q_N(A_s,C_s)(1+R_N(A_s,C_s))^{m-1}\,ds\le K_m.
\tag{5.11b}
\]

There is also a useful exact compensator calculation. During the
mixed-only cycle, entrances to \(D\) raise \(C\) by one and \(D\)-source
exits lower it by one. Since the cycle starts and ends at \(C=0\),

\[
 d_M\mathbb E\int_0^{\sigma_N}(N+C_s)C_s\,ds
 =\mathbb EJ_N.
\tag{5.12}
\]

Applying the generator to \((C)_2\) gives

\[
 d_M\mathbb E\int_0^{\sigma_N}
 (N+C_s)C_s(C_s-1)\,ds
 =\mathbb E\int_0^{\sigma_N}C_s I(A_s)\,ds,
\tag{5.13}
\]

where \(I(A)\) is the total lower-source intensity of reactions entering
\(D\). It satisfies \(I(A)\le K(1+A^2)\le K(1+R_N)^2\). Taking \(m\ge4\)
in (5.11b) and using the term \(NC\) in \(Q_N\) gives

\[
 \mathbb E\int_0^{\sigma_N}C_sI(A_s)\,ds=O(N^{-1}).
\tag{5.14}
\]

This argument uses occupation bounds rather than assigning genealogical
partners to entrances and exits.

Equation (5.13) and (5.14) first give

\[
 \mathbb E\int_0^{\sigma_N}(C_s)_2\,ds=O(N^{-2}).
\tag{5.15}
\]

Since \(C^2=C+(C)_2\), equation (5.12) can now be written as

\[
 \mathbb EJ_N
 =d_MN\mathbb E\int_0^{\sigma_N}C_s\,ds+O(N^{-1}).
\tag{5.16}
\]

The two bounds in (5.8) prove the first estimate in (5.9). This
explicitly rules out a hidden countable-phase tail: the second factorial
occupation is one order smaller than the first, uniformly from the
regeneration phases used below.

\subsection{\texorpdfstring{6. The two-complex pure-\(C\)
linkage}{6. The two-complex pure-C linkage}}

Now restore

\[
 C\mathop{\longrightarrow}^{\alpha}2C,
 \qquad
 2C\mathop{\longrightarrow}^{\beta}C.
\tag{6.1}
\]

The first reaction lowers \(W\) by one and the second raises it by one.
Let \(\widehat\sigma_N\) be the first of (i) the next return to the
minimal \(A,C\)-phase under the full generator and (ii) exit of \(W\)
from \([N/2,3N/2]\). Let \(M_N^-\) and \(M_N^+\) be the numbers of the
two reactions in (6.1), respectively, before this stopping time.

\subsubsection{Lemma 6.1 (interrupted-cycle
estimate)}

There are constants \(0<c<C<\infty\) such that, for all sufficiently
large \(N\),

\[
 \frac cN\le\mathbb EM_N^-\le\frac CN,\qquad
 \mathbb EM_N^+\le\frac C{N^2},
\tag{6.2}
\]

\[
 \mathbb E(M_N^-+M_N^+)^2\le\frac CN,\qquad
 \mathbb E\widehat\sigma_N^2\le C.
\tag{6.3}
\]

To prove the lemma, cut the path at the first mixed-only phase return or
the first pure-\(C\) mark, whichever occurs first, and repeat after a
mark. Uniformly on shells \(n\in[N/2,3N/2]\), (5.11a)--(5.11b) and
(5.15) imply that an attempt started with uniformly bounded phase
moments has a downward-mark hazard \(O(N^{-1})\), an adverse-mark hazard
\(O(N^{-2})\) when it starts from the atom, and a uniformly bounded
second duration moment. For the lower bound, use the fixed feasible
sequence from the end of Lemma 5.1 to enter \(D\) with probability
bounded away from zero. At the resulting bounded phase with \(C=1\), the
probability that the downward clock rings before the \(D\)-exit clock is
at least \(c/N\). Hence the first attempt has a downward mark with
probability at least \(c/N\).

A downward mark changes the phase by one bounded jump. Starting from
that marked phase, the chance of any further mark before mixed return is
\(O(N^{-1})\); in particular, the chance of an adverse mark there is
\(O(N^{-1})\). Thus an adverse mark following a downward mark has
probability \(O(N^{-2})\). A first adverse mark already has probability
\(O(N^{-2})\) by (5.15). Iterating this estimate gives a geometric
interruption expansion: the probability of at least one mark is
\(O(N^{-1})\), while, conditional on a mark, the remaining number of
marks and the remaining return duration have uniformly bounded second
moments. Leaving the shell window would require at least \(N/2\) marks,
so its probability is smaller than every power of \(N^{-1}\). The
high-power phase bounds make its contribution to every first and second
endpoint estimate displayed below smaller than every power as well. This
proves (6.2)--(6.3); it does not assume that the first mixed-only cycle
is uninterrupted or continue a rare escaped path without control.

Since

\[
 \Delta W=M_N^+-M_N^-,
\]

Lemma 6.1 gives constants \(0<d_-<D_-<\infty\) such that

\[
 -\frac{D_-}{N}\le \mathbb E\Delta W
 \le-\frac{d_-}{N},
\qquad
 \mathbb E(\Delta W)^2=O(N^{-1}).
\tag{6.4}
\]

At a minimal phase, \(C=0\), hence \(B=W=N\), and

\[
 U(e_N)=p_Aa_*+p_BN.
\tag{6.5}
\]

On the atom-return event the endpoint has the same \(A,C\)-phase, so its
\(U\)-increment is exactly \(p_B\Delta W\). The shell-exit event
contributes \(o(1)\) to the expected \(\Phi\)-increment by the endpoint
estimate in Lemma 6.1. Thus (6.3)--(6.4) give

\[
 \mathbb E_{e_N}[\Phi(X_{\widehat\sigma_N})-\Phi(e_N)]
 \le-2p_B^2d_-+o(1).
\tag{6.6}
\]

Thus the regeneration drift is bounded above by a negative constant for
all sufficiently large \(N\), while its physical duration and endpoint
second moments are uniformly bounded.

This is the rigorous form of signed service. It does not wait for an
upward record of the autonomous \(C\rightleftarrows2C\) chain. It only
uses the \(O(N^{-1})\)-long visits to \(C=1\) created by the physical
mixed linkage. The factorial record-time obstruction for starting from
large \(C\) is therefore irrelevant; large \(C\) is handled by (4.11).

If the mixed support has no \(0\), the states \(A=C=0\), with \(B\)
fixed, are singleton closed classes. There is one further boundary
family: when \(A=B=0\) and \(C>0\), only the pure linkage is enabled.
For \(\{C,2C\}\) this is the linear-birth/quadratic-death chain on the
positive integers, hence is positive recurrent. Every remaining closed
class uses the positive atom \(a_*=1\), and the preceding argument is
unchanged.

\subsection{\texorpdfstring{7. The three-complex pure-\(C\)
linkage}{7. The three-complex pure-C linkage}}

Suppose instead that \(L_C=\{0,C,2C\}\). Let \(\rho_k\) be the aggregate
rate of the edges \(0\to kC\), for \(k=1,2\). At least one \(\rho_k\) is
positive. Every such reaction lowers \(W\) by \(k\).

The limiting fast trace can be written explicitly. For every lower
target \(j\), put \(q_j=\kappa_{Dj}/d_M\). Start with the effective
one-species generator from (5.3), and add the bounded batch transitions

\[
 a\longmapsto a+j_1+\cdots+j_k
 \quad\hbox{at rate}\quad
 \rho_kq_{j_1}\cdots q_{j_k}.
\tag{7.1}
\]

These transitions represent a \(0\to kC\) event followed by \(k\)
successive \(D\)-exits. They have constant rates and bounded increments.
The original effective generator has linear death in the \(\{0,A\}\)
case and quadratic death in the other three cases. Hence, for every
sufficiently large integer \(m\), the augmented generator
\(\overline{\mathcal L}\) satisfies

\[
 \overline{\mathcal L}(1+a)^m\le K_m-c_m(1+a)^m
\tag{7.2}
\]

outside a finite set. Thus its classwise minimal atom has return times
and endpoint values of every required polynomial moment.

This limiting statement is uniform for the physical chain. Let
\(\widetilde\sigma_N\) be the first of the next minimal-phase return and
exit of \(W\) from \([N/2,3N/2]\). The calculation leading to (5.11a),
with the constant-rate \(0\)-sources and the at-most-linear
\(C\)-sources included, gives the same high-power phase estimate; every
positive linear term is absorbed by \(NC\), \(C^2\), or the stable
lower-phase drift. On a fixed bounded phase set, a newly created \(C\)
particle is cleared at rate at least \(d_MN\). The probability that a
nonzero-source pure reaction interrupts all \(k\) clearings is therefore
\(O(N^{-1})\), uniformly on that set. The high-power estimate makes the
complement uniformly negligible. Combining this with (7.2) and the same
finite-set minorization used in Lemma 5.1 proves

\[
 \mathbb E\widetilde\sigma_N^2\le K.
\tag{7.3}
\]

Leaving the shell window requires at least \(N/4\) pure reactions. The
stopped high-power estimates give moments of every order for their
count, so the shell-exit probability and its first and second endpoint
contributions are smaller than every prescribed power of \(N^{-1}\).

Let \(Z_{k,N}\) count the \(0\to kC\) reactions during this cycle, and
let \(H_N\) count all pure reactions whose source is \(C\) or \(2C\).
The preceding estimates and their stopped high-power versions give

\[
 \mathbb EH_N=O(N^{-1}),\qquad
 \mathbb EH_N^2=O(N^{-1}),\qquad
 \mathbb E\Big(\sum_{k=1}^2 kZ_{k,N}\Big)^2\le K.
\tag{7.4}
\]

The return kernels converge on bounded phases to (7.1), and (7.2)
supplies uniform integrability. Consequently

\[
 \mathbb E\sum_{k=1}^2 kZ_{k,N}=d_0+o(1)
\tag{7.5}
\]

for a constant \(d_0>0\). Strict positivity follows already from the
positive probability that the first reaction at the atom is an enabled
\(0\)-source reaction. Mixed reactions leave \(W\) unchanged. Equations
(7.4)--(7.5), with the bounded signed increments of the \(H_N\)
reactions, therefore give

\[
 \mathbb E\Delta W=-d_0+o(1),\qquad
 \mathbb E(\Delta W)^2\le K.
\tag{7.6}
\]

On the minimal-phase return event the endpoint again has
\(\Delta U=p_B\Delta W\); the shell-exit contribution is \(o(N)\).
Equations (7.3) and (7.6) yield

\[
 \mathbb E[\Phi(X_{\widetilde\sigma_N})-\Phi(e_N)]
 \le-cN
\tag{7.7}
\]

for some \(c>0\) and all large \(N\).

For the architecture (1.5), the lower effective phase is the positive
linear-birth/quadratic-death chain on \(\{A,2A\}\), augmented by the
bounded zero-order batches just described. Its minimal recurrent atom is
\(A=1\). This proves the previously exceptional second architecture
without a finite-\(C\) truncation.

If the mixed support has no \(0\), the boundary \(A=B=0\) carries only
the autonomous pure-\(C\) chain. Its constant/linear births and strictly
negative quadratic \(2C\)-source drift make every closed class positive
recurrent. All other nontrivial classes use the atom \(a_*=1\).

\subsection{8. Hybrid Foster closure}

We now combine the pointwise phase estimate with the regeneration
estimate. Fix a closed irreducible population class. If it is a finite
or singleton boundary class, there is nothing to prove. The autonomous
one-dimensional boundary classes identified at the ends of Sections 6
and 7 are already positive recurrent. In every remaining class choose
the appropriate minimal phase atom from (5.5).

Outside the tube (4.12), stop on first entrance to the tube or to a
fixed finite set. Dynkin's formula and (4.11) charge physical time with
a strict negative margin. On the tube, first hit the minimal phase atom
and then run one regeneration cycle. The high-power bounds above make
the first-hit duration and endpoint moments uniform.

Here is the endpoint check. If the starting phase is \(a>a_*\), then the
preliminary return changes the linear workload by

\[
 p_A(a_*-a)+p_B\Delta W.
\]

The first term is a fixed negative number. In the two-complex case the
positive part of the second term has expectation \(O(N^{-2})\), by the
adverse-mark bound in Lemma 6.1; in the three-complex case its positive
part is \(O(N^{-1})\), by (7.4). Thus the quadratic cross term is
nonpositive for all large \(N\). If \(a=a_*\), the preliminary return is
empty. The appended atom cycle then supplies (6.6) or (7.7).
Consequently there are \(\eta,\delta>0\) and a finite set \(K\) such
that the composite tube stopping time \(\tau_x\) satisfies

\[
 \mathbb E_x\!\left[
 \Phi(X_{\tau_x})-\Phi(x)+\eta\tau_x
 \right]\le-\delta
\tag{8.1}
\]

for every tube state outside \(K\).

Choose \(\eta>0\) smaller than the regeneration drift divided by the
uniform mean cycle duration. Alternating the two stopping rules gives,
up to the first hit of a finite set \(K\),

\[
 \mathbb E\Phi(X_{T_m})
 +\eta\mathbb E T_m
 +\delta\mathbb E N_m
 \le\Phi(X_0),
\tag{8.2}
\]

where \(N_m\) is the number of completed large-\(W\) regeneration cycles
and \(\delta>0\). The endpoint estimates in Lemma 6.1 and Section 7,
together with the high-power bound in Lemma 5.1, justify bounded
optional sampling and passage to the limit. If \(K\) were avoided,
either infinitely much physical time or infinitely many regeneration
cycles would accumulate, contradicting (8.2). Hence \(K\) is reached in
finite mean physical time.

The finite-set trace in a closed irreducible class has finite mean
return, so the class is positive recurrent. This is a physical-time
argument. It does not charge the number of rapid \(D\)-source reactions.

Finally, the process is nonexplosive. Every population-increasing
channel has propensity at most linear: \(0\)-source channels have
constant rates, \(A\)- and \(C\)-source channels have linear rates, and
every quadratic source points to a complex of no larger molecularity.
Comparison of total population with a linear pure-birth process
suffices.

\subsection{9. Exact boundary and audit
conclusions}

The following points are proved above.

\begin{enumerate}
\def\labelenumi{\arabic{enumi}.}
\tightlist
\item
  A linear signed correction to the factorial potential does not close
  the physical-time seam; the bottleneck requires quadratic shell
  growth.
\item
  A rate-dependent positive linear workload \(U\) always exists and
  controls every phase-large state through the proper potential \(U^2\).
\item
  The large positive-\(W\) fast trace has uniform return-time and
  endpoint moments. Its \(C\)-occupation is \(\Theta(W^{-1})\), while
  its \((C)_2\)-occupation is \(O(W^{-2})\).
\item
  For \(C\rightleftarrows2C\), downward signed service is order
  \(W^{-1}\) and adverse service is order \(W^{-2}\); quadratic drift is
  therefore uniformly negative on regeneration cycles.
\item
  When \(0\) is in the pure-\(C\) linkage, downward service has order
  one and the return drift is even stronger.
\item
  Large \(C\) is controlled by quadratic physical death, not by an
  upward record-time policy.
\end{enumerate}

No null or transient closed class occurs in this signed family. The
proof uses the exact supports (1.2)--(1.3); it does not claim a generic
countable-phase theorem for other atlas pairings.

\subsection{10. Exact limit of the superset
interpretation}

The preceding theorem is not monotone under adding complexes to the
pure-\(C\) linkage. The global interface certificate makes the scope gap
exact. Among the 358 residual ordered pairs with one of the four signed
supports in (1.3), 107 available supports contain both \(C\) and \(2C\).
Of those, 69 are already certified by the global tier criterion and 38
remain. The other 161 tier failures do not even contain the two-complex
service support, so no superset version of this note could address them.

The 38 remaining supersets have ten inclusion-minimal ordered pairs:

{\small\def\LTcaptype{none} % do not increment counter
\begin{longtable}[]{@{}ll@{}}
\toprule\noalign{}
signed support & minimal available additions to \(\{C,2C\}\) \\
\midrule\noalign{}
\endhead
\bottomrule\noalign{}
\endlastfoot
\(\{0,2A,B+C\}\) & \(A,\ A+B,\ A+C\) \\
\(\{0,A,2A,B+C\}\) & \(A+B,\ A+C\) \\
\(\{0,A,B+C\}\) & \(2A,\ A+B,\ A+C\) \\
\(\{A,2A,B+C\}\) & \(A+B,\ A+C\) \\
\end{longtable}
}

Thus there are four analytic minimal types: a bounded \(A\)- or
\(2A\)-cofactor, an \(A+B\) source, and an \(A+C\) source. The \(B\)-
and \(2B\)-single-vertex additions are already tier-certified. Adding
\(0\) alone is the exact three-complex case of Section 7.

Each of the four types changes a load-bearing estimate above. An \(A\)-
or \(2A\)-source can fire with \(C=0\), so \(W\) is no longer fixed
during the lower-phase return. An \(A+B\)-source has rate of order \(N\)
as soon as one \(A\) particle is present, and becomes a second fast
clock not contained in the stochastic complement (5.3). An \(A+C\)
intermediate introduces a second reaction that can have to beat the
order-\(N\) \((B+C)\)-clearing clock.

For example, consider an orientation containing

\[
 C\longrightarrow A+C\longrightarrow2C.
\tag{10.1}
\]

At the phase \(A=0,C=1,B=N+1\), the first arrow beats clearing with
probability \(\Theta(N^{-1})\). Conditional on that arrow, the second
again beats clearing with probability \(\Theta(N^{-1})\). Thus this
particular direct service path has probability \(\Theta(N^{-2})\), not
the one-clock probability used in Lemma 6.1.

This observation alone is not a counterexample to an \(N^{-1}\)
\emph{averaged} service rate. The irreducible lower phase has a
positive, though possibly very small, occupation of \(A>0\); on those
visits \(A+C\to2C\) needs to win only one order-\(N\) race. Establishing
the resulting coefficient requires a new countable-phase occupation and
interruption theorem that includes the added reactions. It is not
supplied by Lemmas 5.1 or 6.1.

Accordingly, (10.1) is a countersequence to a direct one-path
restoration argument, while the \(A\), \(2A\), and \(A+B\) cases
invalidate the fixed-\(W\), delete-only-\(D\) trace itself. None is
asserted to be a null-recurrent or transient reaction network. No
deletion or restoration monotonicity is used in the theorem above.
